% Build from docs/guide/: see README.md (latexmk, output in build/).
\documentclass[11pt,a4paper,oneside,openany]{book}
% Shared by handbook.tex and any private handbook that % Shared by handbook.tex and any private handbook that % Shared by handbook.tex and any private handbook that \input{preamble}s it.
% The main file sets \documentclass, the title page, \hypersetup{pdftitle=...}
% and \addbibresource before or after this file; it may define \guidehead first.
\usepackage[T1]{fontenc}
\usepackage[utf8]{inputenc}
\usepackage[USenglish]{babel}
\usepackage{lmodern}
\usepackage[a4paper,margin=2.3cm,includehead,includefoot,headheight=14pt]{geometry}
\usepackage{microtype,setspace,amsmath,amssymb}
\usepackage{placeins}
\usepackage{booktabs,longtable,tabularx,array}
\usepackage[table]{xcolor}
\usepackage{enumitem,needspace,listings,float,graphicx}
\usepackage{tikz}
\usetikzlibrary{arrows.meta,positioning,calc}
\usepackage{fancyhdr,csquotes}
\usepackage[backend=biber,style=alphabetic,sorting=nyt,backref=true,
  giveninits=true,maxbibnames=8]{biblatex}
\setcounter{biburlnumpenalty}{9000}
\setcounter{biburlucpenalty}{9000}
\AtBeginBibliography{\raggedright}
\renewcommand*{\bibfont}{\small}
\usepackage{siunitx}
\sisetup{group-separator={,},group-minimum-digits=4}
\sisetup{retain-explicit-plus=true,group-digits=integer}
\usepackage{xurl}
\usepackage{hyperref}
\definecolor{ctpText}{HTML}{4C4F69}
\definecolor{ctpSurface0}{HTML}{CCD0DA}
\definecolor{ctpBlue}{HTML}{1E66F5}
\definecolor{ctpMauve}{HTML}{8839EF}
\definecolor{ctpGreen}{HTML}{40A02B}
\definecolor{ctpYellow}{HTML}{DF8E1D}
\definecolor{ctpPeach}{HTML}{FE640B}
\definecolor{ctpRed}{HTML}{D20F39}
\hypersetup{colorlinks=true,linkcolor=ctpBlue,citecolor=ctpBlue,urlcolor=ctpBlue}
\setstretch{1.15}
\widowpenalty=10000
\clubpenalty=10000
\setlength{\parindent}{0pt}
\setlength{\parskip}{0.55em}
\setlist{itemsep=0.25em,topsep=0.4em,leftmargin=1.6em}
\renewcommand{\arraystretch}{1.16}
\setcounter{tocdepth}{1}
\newcolumntype{Y}{>{\raggedright\arraybackslash}X}
\newcolumntype{L}[1]{>{\raggedright\arraybackslash}p{#1}}
\newcommand{\code}[1]{\nolinkurl{#1}}
\newcommand{\contract}{\textcolor{ctpBlue}{\textbf{[C] Contract}}}
\newcommand{\evidence}{\textcolor{ctpGreen}{\textbf{[E] Evidence}}}
\newcommand{\gap}{\textcolor{ctpPeach}{\textbf{[G] Gap}}}
\newcommand{\review}{\textcolor{ctpMauve}{\textbf{[R] Review}}}
\lstset{basicstyle=\ttfamily\footnotesize,backgroundcolor=\color{ctpSurface0!25},
  frame=single,rulecolor=\color{ctpText!35},breaklines=true,columns=fullflexible,
  keepspaces=true,showstringspaces=false}
\tikzset{flow/.style={-Latex,thick,draw=ctpText},
  box/.style={rounded corners=2mm,thick,align=center,inner sep=6pt,font=\small},
  work/.style={box,draw=ctpBlue,fill=ctpBlue!10},
  check/.style={box,draw=ctpMauve,fill=ctpMauve!10},
  result/.style={box,draw=ctpGreen,fill=ctpGreen!10},
  open/.style={box,draw=ctpPeach,fill=ctpPeach!10}}
\providecommand{\guidehead}{Agent workflows}
\pagestyle{fancy}
\fancyhf{}
\fancyhead[L]{\nouppercase{\leftmark}}
\fancyhead[R]{\guidehead}
\fancyfoot[C]{\thepage}
\renewcommand{\chaptermark}[1]{\markboth{\thechapter.\ #1}{}}
% Insertion point: a private handbook defines \csname hook@<name>\endcsname to
% add its own sections here; the public build leaves every hook empty.
\newcommand{\hook}[1]{\ifcsname hook@#1\endcsname\csname hook@#1\endcsname\fi}
s it.
% The main file sets \documentclass, the title page, \hypersetup{pdftitle=...}
% and \addbibresource before or after this file; it may define \guidehead first.
\usepackage[T1]{fontenc}
\usepackage[utf8]{inputenc}
\usepackage[USenglish]{babel}
\usepackage{lmodern}
\usepackage[a4paper,margin=2.3cm,includehead,includefoot,headheight=14pt]{geometry}
\usepackage{microtype,setspace,amsmath,amssymb}
\usepackage{placeins}
\usepackage{booktabs,longtable,tabularx,array}
\usepackage[table]{xcolor}
\usepackage{enumitem,needspace,listings,float,graphicx}
\usepackage{tikz}
\usetikzlibrary{arrows.meta,positioning,calc}
\usepackage{fancyhdr,csquotes}
\usepackage[backend=biber,style=alphabetic,sorting=nyt,backref=true,
  giveninits=true,maxbibnames=8]{biblatex}
\setcounter{biburlnumpenalty}{9000}
\setcounter{biburlucpenalty}{9000}
\AtBeginBibliography{\raggedright}
\renewcommand*{\bibfont}{\small}
\usepackage{siunitx}
\sisetup{group-separator={,},group-minimum-digits=4}
\sisetup{retain-explicit-plus=true,group-digits=integer}
\usepackage{xurl}
\usepackage{hyperref}
\definecolor{ctpText}{HTML}{4C4F69}
\definecolor{ctpSurface0}{HTML}{CCD0DA}
\definecolor{ctpBlue}{HTML}{1E66F5}
\definecolor{ctpMauve}{HTML}{8839EF}
\definecolor{ctpGreen}{HTML}{40A02B}
\definecolor{ctpYellow}{HTML}{DF8E1D}
\definecolor{ctpPeach}{HTML}{FE640B}
\definecolor{ctpRed}{HTML}{D20F39}
\hypersetup{colorlinks=true,linkcolor=ctpBlue,citecolor=ctpBlue,urlcolor=ctpBlue}
\setstretch{1.15}
\widowpenalty=10000
\clubpenalty=10000
\setlength{\parindent}{0pt}
\setlength{\parskip}{0.55em}
\setlist{itemsep=0.25em,topsep=0.4em,leftmargin=1.6em}
\renewcommand{\arraystretch}{1.16}
\setcounter{tocdepth}{1}
\newcolumntype{Y}{>{\raggedright\arraybackslash}X}
\newcolumntype{L}[1]{>{\raggedright\arraybackslash}p{#1}}
\newcommand{\code}[1]{\nolinkurl{#1}}
\newcommand{\contract}{\textcolor{ctpBlue}{\textbf{[C] Contract}}}
\newcommand{\evidence}{\textcolor{ctpGreen}{\textbf{[E] Evidence}}}
\newcommand{\gap}{\textcolor{ctpPeach}{\textbf{[G] Gap}}}
\newcommand{\review}{\textcolor{ctpMauve}{\textbf{[R] Review}}}
\lstset{basicstyle=\ttfamily\footnotesize,backgroundcolor=\color{ctpSurface0!25},
  frame=single,rulecolor=\color{ctpText!35},breaklines=true,columns=fullflexible,
  keepspaces=true,showstringspaces=false}
\tikzset{flow/.style={-Latex,thick,draw=ctpText},
  box/.style={rounded corners=2mm,thick,align=center,inner sep=6pt,font=\small},
  work/.style={box,draw=ctpBlue,fill=ctpBlue!10},
  check/.style={box,draw=ctpMauve,fill=ctpMauve!10},
  result/.style={box,draw=ctpGreen,fill=ctpGreen!10},
  open/.style={box,draw=ctpPeach,fill=ctpPeach!10}}
\providecommand{\guidehead}{Agent workflows}
\pagestyle{fancy}
\fancyhf{}
\fancyhead[L]{\nouppercase{\leftmark}}
\fancyhead[R]{\guidehead}
\fancyfoot[C]{\thepage}
\renewcommand{\chaptermark}[1]{\markboth{\thechapter.\ #1}{}}
% Insertion point: a private handbook defines \csname hook@<name>\endcsname to
% add its own sections here; the public build leaves every hook empty.
\newcommand{\hook}[1]{\ifcsname hook@#1\endcsname\csname hook@#1\endcsname\fi}
s it.
% The main file sets \documentclass, the title page, \hypersetup{pdftitle=...}
% and \addbibresource before or after this file; it may define \guidehead first.
\usepackage[T1]{fontenc}
\usepackage[utf8]{inputenc}
\usepackage[USenglish]{babel}
\usepackage{lmodern}
\usepackage[a4paper,margin=2.3cm,includehead,includefoot,headheight=14pt]{geometry}
\usepackage{microtype,setspace,amsmath,amssymb}
\usepackage{placeins}
\usepackage{booktabs,longtable,tabularx,array}
\usepackage[table]{xcolor}
\usepackage{enumitem,needspace,listings,float,graphicx}
\usepackage{tikz}
\usetikzlibrary{arrows.meta,positioning,calc}
\usepackage{fancyhdr,csquotes}
\usepackage[backend=biber,style=alphabetic,sorting=nyt,backref=true,
  giveninits=true,maxbibnames=8]{biblatex}
\setcounter{biburlnumpenalty}{9000}
\setcounter{biburlucpenalty}{9000}
\AtBeginBibliography{\raggedright}
\renewcommand*{\bibfont}{\small}
\usepackage{siunitx}
\sisetup{group-separator={,},group-minimum-digits=4}
\sisetup{retain-explicit-plus=true,group-digits=integer}
\usepackage{xurl}
\usepackage{hyperref}
\definecolor{ctpText}{HTML}{4C4F69}
\definecolor{ctpSurface0}{HTML}{CCD0DA}
\definecolor{ctpBlue}{HTML}{1E66F5}
\definecolor{ctpMauve}{HTML}{8839EF}
\definecolor{ctpGreen}{HTML}{40A02B}
\definecolor{ctpYellow}{HTML}{DF8E1D}
\definecolor{ctpPeach}{HTML}{FE640B}
\definecolor{ctpRed}{HTML}{D20F39}
\hypersetup{colorlinks=true,linkcolor=ctpBlue,citecolor=ctpBlue,urlcolor=ctpBlue}
\setstretch{1.15}
\widowpenalty=10000
\clubpenalty=10000
\setlength{\parindent}{0pt}
\setlength{\parskip}{0.55em}
\setlist{itemsep=0.25em,topsep=0.4em,leftmargin=1.6em}
\renewcommand{\arraystretch}{1.16}
\setcounter{tocdepth}{1}
\newcolumntype{Y}{>{\raggedright\arraybackslash}X}
\newcolumntype{L}[1]{>{\raggedright\arraybackslash}p{#1}}
\newcommand{\code}[1]{\nolinkurl{#1}}
\newcommand{\contract}{\textcolor{ctpBlue}{\textbf{[C] Contract}}}
\newcommand{\evidence}{\textcolor{ctpGreen}{\textbf{[E] Evidence}}}
\newcommand{\gap}{\textcolor{ctpPeach}{\textbf{[G] Gap}}}
\newcommand{\review}{\textcolor{ctpMauve}{\textbf{[R] Review}}}
\lstset{basicstyle=\ttfamily\footnotesize,backgroundcolor=\color{ctpSurface0!25},
  frame=single,rulecolor=\color{ctpText!35},breaklines=true,columns=fullflexible,
  keepspaces=true,showstringspaces=false}
\tikzset{flow/.style={-Latex,thick,draw=ctpText},
  box/.style={rounded corners=2mm,thick,align=center,inner sep=6pt,font=\small},
  work/.style={box,draw=ctpBlue,fill=ctpBlue!10},
  check/.style={box,draw=ctpMauve,fill=ctpMauve!10},
  result/.style={box,draw=ctpGreen,fill=ctpGreen!10},
  open/.style={box,draw=ctpPeach,fill=ctpPeach!10}}
\providecommand{\guidehead}{Agent workflows}
\pagestyle{fancy}
\fancyhf{}
\fancyhead[L]{\nouppercase{\leftmark}}
\fancyhead[R]{\guidehead}
\fancyfoot[C]{\thepage}
\renewcommand{\chaptermark}[1]{\markboth{\thechapter.\ #1}{}}
% Insertion point: a private handbook defines \csname hook@<name>\endcsname to
% add its own sections here; the public build leaves every hook empty.
\newcommand{\hook}[1]{\ifcsname hook@#1\endcsname\csname hook@#1\endcsname\fi}

\hypersetup{pdftitle={Working with Agents on Scientific Projects},pdfauthor={Scientific agents}}
\title{\textbf{Working with Agents on Scientific Projects}\\[0.4em]
 \large Daily work, scientific code, and reports you can check}
\author{Jairo Tavizon}
\date{Handbook for the skills of the repository, October 2026}
\addbibresource{bib/handbook.bib}
\begin{document}
\frontmatter
\pagestyle{plain}
\hypersetup{pageanchor=false}
\maketitle
\hypersetup{pageanchor=true}
\chapter*{Purpose and reading routes}
\addcontentsline{toc}{chapter}{Purpose and reading routes}

\begin{headline}
  This handbook shows how to turn a research question into work its owner can understand,
  check and continue when agents do part of it.
\end{headline}

It starts with an ordinary working day, develops the scientific-coding and report-writing
flows in detail, then connects them to remote execution, coordinated changes and
recovery. Each flow explains the inputs, decisions, intermediate evidence and conditions
for handing work back.

It is for people who use Git and write scientific code but are new to coordinating
agents. The agents read the skills of the repository, not this book; the book names each
skill where it applies, and every claim cites the skill file and section behind it. Take
what helps and leave the rest.

\begin{table}[H]
  \small
  \begin{tabularx}{\textwidth}{L{0.3\textwidth}Y}
    \toprule
    When you are & Read \\
    \midrule
    Starting a day & Chapter~\ref{ch:daily}: pick one task and its flow. \\
    Producing a result & Chapter~\ref{ch:code} for scientific code,
    Chapter~\ref{ch:reports} for a report; Chapter~\ref{ch:example} walks one small task
    from question to reviewed report. \\
    At a transition & Chapters~\ref{ch:remote}--\ref{ch:recovery}: remote results, another
    writer, an interruption, preservation. \\
    Adapting the skills & Chapter~\ref{ch:yours}. \\
    \bottomrule
  \end{tabularx}
  \caption{Reading routes.}
  \label{tab:routes}
\end{table}

Skill names in the text refer to the folders under \code{skills/} in the repository
\url{https://github.com/jtavizon98/scientific-agents}; its README lists them all.

\paragraph{Citations.} Source handling follows one method: readable bibliography labels,
claim-level pinpoints, and links back from a source entry to the pages that cite it. A
citation such as \cite[Validation order]{skill:implementation}
names the owning source and the section to inspect. Most workflow sources have no printed
pages, so their headings are the locators; Chapter~\ref{ch:reports} applies the same method
to papers and figures. To check a claim:
\begin{enumerate}
  \item follow the citation to its entry in the sources;
  \item open the named source and inspect the pinpoint;
  \item to find other uses of the source, follow the entry's return links.
\end{enumerate}
The bibliography locates evidence; it does not prove that an uninspected source supports
a claim.

\begingroup
\setstretch{1.0}
\tableofcontents
\endgroup

\chapter*{Essential terms}
\addcontentsline{toc}{chapter}{Essential terms}
\begin{description}
  \item[Control surface] Where the owner interacts: for example, a phone or a laptop client. It
        need not be the machine executing the work.
  \item[Runtime and writing host] The machine running the agent and hosting the local checkout it
        owns. A separate execution host may run an analysis job.
  \item[Session] A bounded task with an owner, current state, and intended result. A client
        conversation may carry it, but durable records preserve it.
  \item[Worktree] A separate Git checkout used to isolate a writer's files. Each writing boundary
        has an explicit owner.
  \item[Contract] The agreed question, scope, deliverables, validation, and applicable review and
        handoff conditions.
  \item[Oracle] An independent expected result or reference used to judge an implementation. The
        reference's authority and limitations must be known.
  \item[Pinpoint] The exact page, section, equation, figure, table, or record supporting a
        citation at its point of use.
  \item[Handback] Evidence returned by a worker or monitor. The parent still owns any remaining
        analysis, review, and delivery obligations.
  \item[Review gate] A named check required before a claim or artifact can be accepted. Copy,
        source, domain, and rendered-page checks answer different questions.
  \item[Human review] The owner's judgment of the delivered artifact or scientific decision. A
        report's PDF acceptance is distinct from its technical checks.
\end{description}

\section*{Reading the diagrams and labels}

Diagrams use four colours, each with one meaning; the text inside a box always names the
step, so colour is a second cue, never the only one.

\begin{center}
  \begin{tikzpicture}[node distance=4mm]
    \node[work] (w) {Work an agent does};
    \node[check,right=of w] (c) {A check or review};
    \node[result,right=of c] (r) {A result};
    \node[open,right=of r] (o) {A step for a person};
  \end{tikzpicture}
\end{center}

Boxes in the text have fixed roles:
\begin{itemize}
  \item a grey \textbf{headline} bar opens a section with its point;
  \item an orange \textbf{choice} box lists the options for a decision, with a recommendation;
  \item a red \textbf{rule} box states a limit that holds without exception.
\end{itemize}
In the worked examples, \contract, \evidence, \gap\ and \review\ mark requirements,
observations, missing evidence and checks.


\mainmatter
\pagestyle{fancy}
\chapter{Start and Continue Daily Work}\label{ch:daily}
\section{Choose the result before the client}

A useful starting request names what should be different when the work is done. For
example, ``explain where these two cutflows first diverge'' defines an investigation;
``change the selection and show its effect'' additionally defines implementation and
scientific validation. Before acting, identify the owning repository or document, the
evidence available, and the result the reader will inspect. This determines whether the
work needs a short interactive exchange or a durable session \cite[rules
  1--3]{skill:coordination}.

For any task beyond a quick edit, start the agent in \code{scientific-mode}. It states
the goal back in its own words, picks the playbook that fits (a study or comparison,
figures and reports, or picking up, pausing and closing work) and copies the playbook's
steps into its todo list, so a skipped step is visible rather than forgotten \cite[Start;
  Playbooks]{skill:mode}.

Routine, mechanical, local work can remain lightweight. A spelling correction or a
bounded inspection does not need the full scaffolding of a remote study. Delegated,
unattended, cross-host, consequential, long-running, substantial, or recovery-worthy work
requires \code{BRIEF.md}, \code{STATE.md}, and \code{RESULT.md}. Add a machine-readable
\code{HANDOFF.json} when programmatic dependencies or cross-harness recovery require one
\cites[rule 3]{skill:coordination}[Handoff marker]{skill:jobs}.

The brief preserves the commitment: question, owner, write boundary, evidence,
verification, allowed actions, and stopping conditions. State records the actual
checkpoint, remaining work, and next safe action. The result states what was produced and
which claims remain limited. These records serve different points in the same task;
keeping state current makes a later handoff much easier \cite[Dispatch
  packet]{skill:delegation}.

For a substantial task, make the completion sentence concrete. ``The comparison report is
ready for my review when the inputs and cutflows have been checked, the report passes
internal reviews, and the exact PDF is available locally'' identifies an assessable
intermediate outcome. Human acceptance remains open. A code-only task can instead name
its implementation, checks, and reading guide; it acquires a report requirement only when
the task contract calls for one \cites[rule 4]{skill:coordination}[Verify and hand
  back]{skill:coding}.

\section{Identify the owner and the execution boundary}

Before choosing commands, distinguish four roles: the control surface, the agent runtime,
the owned checkout, and any execution target. Figure~\ref{fig:roles} shows their
relationship. A phone can steer an agent on a workstation whose local worktree contains
the source while an authorized job runs on another host. The phone's location does not
relocate that source or change the job's owner \cite[GUI, terminal and cross-client
  sessions]{skill:workflow-gui}.

\begin{figure}[htbp]
  \centering
  \begin{tikzpicture}[node distance=10mm]
    \node[work,text width=110mm] (ui) {\textbf{Control surface}\\phone or laptop;
      choose the intended session};
    \node[work,text width=110mm,below=of ui] (runtime) {\textbf{Agent runtime and owned worktree}\\
      observed host, repository, branch, writing owner};
    \node[check,text width=110mm,below=of runtime] (run) {\textbf{Execution target, when needed}\\
      capable environment and explicitly authorized job};
    \draw[flow] (ui)--node[right,font=\footnotesize]{steer}(runtime);
    \draw[flow] (runtime)--node[right,font=\footnotesize]{run / inspect}(run);
  \end{tikzpicture}
  \caption{Control, source ownership, and execution are separate roles. This is a
    routing model, not evidence of a particular live connection.}
  \label{fig:roles}
\end{figure}

Inspect the literal host, Git state, worktree, and active session before editing. An
unfamiliar dirty path needs ownership evidence. Preserve it while resolving the
uncertainty; do not reset it to make the checkout look convenient. On the same machine,
use an isolated worktree when report or software edits would overlap another task
\cite[Start a writing task; Commit and hand off]{skill:git}.

\hook{daily-continuation}

\section{Keep the work loop small and the ending explicit}

Within the owned boundary, alternate a coherent change with the evidence needed to assess
it. Remote execution belongs inside that loop whenever the check needs remote data or
software. Figure~\ref{fig:loop} shows the dependency without forcing every task through a
report or putting all validation before execution.

\begin{figure}[htbp]
  \centering
  \begin{tikzpicture}[node distance=14mm]
    \node[work,text width=45mm] (do) {\textbf{Implement or investigate}\\
      bounded question and owned change};
    \node[check,text width=45mm,right=18mm of do] (test) {\textbf{Run and inspect}\\
      local or authorized remote evidence};
    \node[result,text width=108mm,below=22mm of do.south east,anchor=north] (hand)
    {\textbf{Hand back the promised result}\\applicable reviews, artifact,
      remaining decisions, and publication state};
    \draw[flow] (do)--(test);
    \draw[flow] (test.north) to[bend right=35] node[above,font=\footnotesize]
    {failed check: revise} (do.north);
    \draw[flow] (test.south)--(hand.north east);
  \end{tikzpicture}
  \caption{A shared work loop with conditional outputs. Source publication can
    occur at authorized durability checkpoints; report review can return to revision
    \cites[Verify and hand back]{skill:coding}[Review and release]{skill:report-production}
    [Commit and hand off]{skill:git}.}
  \label{fig:loop}
\end{figure}

Close only the gates named by the task. Code checks, source publication, PDF
availability, external delivery, notification, and human acceptance are separately
observable states. Record an absent requirement as not applicable and an unmet
requirement as pending. A substantial report remains \texttt{awaiting human review} after
internal checks and requested delivery until its owner accepts, supersedes, or closes it.
A version that an agent replaces before the owner has reviewed it is marked
\texttt{replaced before review}; \texttt{superseded} is reserved for the owner's own
review or decision \cite[Review and release]{skill:report-production}.

\chapter{Investigate and Settle Decisions}\label{ch:decisions}
\section{Build an evidence base that answers the question}

\begin{headline}
  Start from the source that owns the behavior, and keep rules, observations, inferences
  and proposals apart.
\end{headline}

\begin{table}[H]
  \small
  \begin{tabularx}{\textwidth}{L{0.25\textwidth}Y}
    \toprule
    For & Inspect \\
    \midrule
    An analysis mismatch & The relevant definitions, implementation, input identity and
    output evidence. \\
    A workflow change & Its canonical procedure and current use. \\
    \bottomrule
  \end{tabularx}
  \caption{Where an investigation starts.}
  \label{tab:decisions-start}
\end{table}

\begin{rulebox}[Search safety]
  Search only explicit safe roots. Inspect known files and narrow remote paths; never scan
  the whole workspace or a mounted remote filesystem recursively \cite[Remote checkouts and
    execution]{skill:remote}.
\end{rulebox}

\begin{table}[H]
  \small
  \begin{tabularx}{\textwidth}{L{0.2\textwidth}Y}
    \toprule
    Statement & Says \\
    \midrule
    Rule & What the owning source requires. \\
    Observation & What a check actually found. \\
    Inference & What explains the observations. \\
    Proposal & A change not yet accepted. \\
    \bottomrule
  \end{tabularx}
  \caption{Four kinds of statement. For example, a successful run in one client is
    an observation supporting a client preference; it
    does not show that an unattended notification mechanism works \cite[Source traceability
      and corroboration]{skill:report-production}.}
  \label{tab:decisions-statements}
\end{table}

When evidence disagrees, find the first disagreement and the authority of each side:
\begin{itemize}
  \item a session can show that a run used a configuration; it cannot redefine the analysis
        convention;
  \item a policy file can require validation; it cannot prove that a run passed;
  \item a conclusion that depends on both cites both \cite[Source traceability and
          corroboration]{skill:report-production}.
\end{itemize}

\section{Use the planning pass for material choices}

\begin{headline}
  Before a material change, settle the owner's decisions in dependency order and confirm
  them as one record.
\end{headline}

The grilling pass applies before workflow or policy changes, architecture or
cross-repository work, physics or analysis design, and substantial report revisions
(Figure~\ref{fig:grilling}) \cite[Establish the decision tree; Work the
  frontier]{skill:grilling}.

\begin{figure}[H]
  \centering
  \begin{tikzpicture}[node distance=5mm and 6mm]
    \node[work,text width=30mm] (facts) {Inspect the facts};
    \node[work,text width=30mm,right=of facts] (split) {Separate owner choices from
      facts};
    \node[work,text width=30mm,right=of split] (order) {Order choices by dependency};
    \node[open,text width=30mm,below=8mm of order] (ask) {Ask the choices whose
      prerequisites are known};
    \node[open,text width=30mm,left=of ask] (confirm) {Confirm the decision record};
    \node[work,text width=30mm,left=of confirm] (impl) {Implement};
    \draw[flow] (facts)--(split);
    \draw[flow] (split)--(order);
    \draw[flow] (order)--(ask);
    \draw[back] (ask.east) to[bend right=50] node[right,lbl]{next\\round} (order.east);
    \draw[flow] (ask)--(confirm);
    \draw[flow] (confirm)--(impl);
  \end{tikzpicture}
  \caption{The grilling pass. Ordinary naming and formatting follow repository conventions
    without another round.}
  \label{fig:grilling}
\end{figure}

For example, a request to make an analysis script easier to understand first needs to
know whether it must preserve the scientific definition or implement a new one: that
decides the validation oracle. Then come the owning layer, the bounded input slice and
the acceptance evidence \cite[Establish the scientific contract]{skill:coding}.

The decision record keeps:
\begin{itemize}
  \item settled scope, non-goals and assumptions, with their evidence;
  \item affected owners and the implementation order;
  \item checks, the recovery path and remaining blockers.
\end{itemize}
Confirm it before implementation where the gate applies. If the owner already specified
every material choice, record them directly; do not manufacture a new round \cite[Confirm
  the decision record; Work the frontier]{skill:grilling}.

\section{Carry authorization with the action}

\begin{headline}
  Confirming a plan establishes understanding. A consequential action also needs an exact
  authorization boundary.
\end{headline}

Table~\ref{tab:consequence} is a compact guide; the consequential-actions procedure has
the full rules.

\begin{table}[H]
  \small
  \begin{tabularx}{\textwidth}{L{0.25\textwidth}Y}
    \toprule
    Action & What must be established \\
    \midrule
    Scoped inspection & Relevant target and safe read boundary. \\
    \rowcolor{ctpSurface0!30}
    Local reversible work & Assigned writing boundary and task scope. \\
    Own branch and review inbox & Nothing beyond the task: push the task's own branch or
    deliver a ready artifact to the reader's review inbox, then report it. \\
    \rowcolor{ctpSurface0!30}
    Shared change & Exact intent, target, state, recovery, and verification. Includes
    push to a main or shared branch, remote checkout update, and external PDF delivery. \\
    Operational action & Explicit task authorization for the job, service, mount,
    or other operation; credential initialization remains user-owned where required. \\
    \rowcolor{ctpSurface0!30}
    Destructive action & Specific authorization and recovery evidence for the exact
    target; cleanup additionally binds the immutable plan and batch. \\
    \bottomrule
  \end{tabularx}
  \caption{Action classes and their prerequisites
    \cite[Preflight]{skill:actions}.}
  \label{tab:consequence}
\end{table}

\begin{itemize}
  \item Record who authorized what, from which instruction, for which target and scope.
  \item The atomic steps of an authorized action need no new permission question each.
  \item Recheck when the target, state, scope, consequence or instruction changes.
  \item A successful inspection or a convenient recovery path never supplies missing intent
        \cite[Preflight]{skill:actions}.
\end{itemize}

\section{Decide, record, report}

\begin{headline}
  Between the gates, an agent does reversible work without asking and says what it did
  and why.
  A blocking scientific choice follows Figure~\ref{fig:drr}.
\end{headline}

\begin{figure}[H]
  \centering
  \begin{tikzpicture}[node distance=6mm and 8mm]
    \node[check,text width=42mm] (q) {A scientific choice blocks progress\\(a definition,
      which comparison answers the question)};
    \node[check,text width=34mm,right=12mm of q] (cheap) {Can a cheap run settle it?};
    \node[work,text width=34mm,right=12mm of cheap] (run) {Run it and decide from the
      result};
    \node[check,text width=34mm,below=10mm of cheap] (def) {Is a default defensible, and
      the purpose unchanged?};
    \node[work,text width=34mm,right=12mm of def] (rec) {Take the default best supported by the
      sources and what is already agreed, record it, keep working};
    \node[open,text width=34mm,below=10mm of def] (ask) {Ask the user before acting};
    \draw[flow] (q)--(cheap);
    \draw[flow] (cheap)--node[above,lbl]{yes}(run);
    \draw[flow] (cheap)--node[right,lbl]{no}(def);
    \draw[flow] (def)--node[above,lbl]{yes}(rec);
    \draw[flow] (def)--node[right,lbl]{no}(ask);
  \end{tikzpicture}
  \caption{Decide, record, report \cites{skill:prototype}[Decide, record,
      report]{skill:actions}.}
  \label{fig:drr}
\end{figure}

A recorded default goes into the session's table of agent defaults with the alternative,
the reason, the evidence and how to reverse it. Every result that depends on an
unreviewed default stays awaiting human review, and the next report lists those defaults
first. The work keeps moving while the scientific judgment stays visible and reversible.

\begin{rulebox}[Never a default]
  A confirmed decision record or an agreed convention is never a default to override.
\end{rulebox}

\chapter{Develop and Validate Scientific Code}\label{ch:code}
\section{From a physical transformation to a stage record}

\begin{headline}
  Scientific code exposes the transformation we want to understand. Write its stage
  record before editing: it reveals missing decisions early, and a blank source or oracle is
  a real gap.
\end{headline}

A reader should find the entry point, follow physically named intermediate states, and
locate the selection, weight, formula or convention behind an output. The workflow
applies to substantial scientific code, physics behavior, architecture and explanatory
rewrites; a framework's native conventions stay authoritative for its interfaces
\cite[Scientific coding; Expose the procedure]{skill:coding}.

\paragraph{Worked example.} Reorganize a script that reads an approved nominal input,
applies an approved selection and fills a weighted histogram of an existing observable, so
that it reads clearly while behaving the same. The example fixes no cuts or bin values and
reports no measured result; its scientific choices must come from the owning analysis source.
Table~\ref{tab:stage} is its record.

\begingroup
\small
\begin{longtable}{L{0.27\textwidth}L{0.66\textwidth}}
  \caption{Applying the required stage record to the illustrative histogram rewrite
    \cite[Establish the scientific contract]{skill:coding}.}\label{tab:stage}\\
  \toprule
  Field & Record before editing \\
  \midrule
  \endfirsthead
  \toprule Field & Record before editing \\ \midrule
  \endhead
  Question & Can the existing transformation be made readable while preserving
  its declared behavior? \\
  \rowcolor{ctpSurface0!30}
  Authoritative source & Exact analysis definition and repository section owning
  selection, observable, weighting, and binning. \\
  Owning layer & Repository and module that own analysis combinations; framework
  calibration stays in its framework layer. \\
  \rowcolor{ctpSurface0!30}
  Input and output & Frozen input identity and schema; histogram artifact plus
  event accounting and comparison evidence. \\
  Named stages & Load inputs; select objects/events; compute observable; apply
  weights; fill and persist the result. Use the actual analysis order. \\
  \rowcolor{ctpSurface0!30}
  Invariants and conventions & Units, selection boundaries, object ordering,
  invalid-value treatment, weight signs, bin edges, and flow handling. \\
  Scale and performance & Expected input size; baseline host, CPU allocation,
  I/O location, wall time, and peak memory where performance matters. \\
  \rowcolor{ctpSurface0!30}
  Oracle & Independent small examples and an inspected baseline artifact whose
  definition is understood; declare comparison tolerances. \\
  Acceptance evidence & Stage comparisons, intended-input execution, focused
  invariant checks, required reviews, and measured performance if applicable. \\
  \rowcolor{ctpSurface0!30}
  Non-goals & No new observable, changed cut, changed normalization, or new
  execution framework unless separately agreed. \\
  \bottomrule
\end{longtable}
\endgroup

\contract\ The transformation and its boundary are now explicit. A blank authoritative source or
oracle would mean we do not yet know what agreement would demonstrate.
\gap\ If the existing output is known to contain a bug, preserving it is not a valid
oracle for that behavior: confirm the intended correction and change the expected result.
A refactor and a definition change can share infrastructure, but their acceptance
questions differ \cite[Numerical evidence]{skill:implementation}.

\section{Give each definition one owner}

\begin{headline}
  Trace the output backward and record one implementation owner for each definition, so a
  convenient copy never becomes a second canonical analysis.
\end{headline}

For the example, find where the binning is chosen, the observable calculated, the event
weight assembled and the selection applied. Independent diagnostics may still check the
production result when their role is explicit (Table~\ref{tab:code-owners})
\cite[Ownership and framework composition]{skill:implementation}.

\begin{table}[H]
  \small
  \begin{tabularx}{\textwidth}{L{0.36\textwidth}Y}
    \toprule
    Definition & Owner \\
    \midrule
    Official calibration and framework decorations & The framework layer. \\
    \rowcolor{ctpSurface0!30}
    Combinations of stable objects & The analysis layer. \\
    Mathematics shared by reconstructed and particle-level definitions & A shared helper,
    not copied formulas. The two definitions may differ intentionally, and the helper must not
    erase a meaningful difference between them. \\
    \rowcolor{ctpSurface0!30}
    Host paths & Runtime configuration, kept apart from physics choices where supported. \\
    \bottomrule
  \end{tabularx}
  \caption{Where definitions live \cite[Ownership and framework
      composition]{skill:implementation}.}
  \label{tab:code-owners}
\end{table}

\paragraph{Read in scientific order.} In the example one should see which selected
events enter the observable and which weight reaches the fill. Small pure helpers make an
independently specified operation easy to test; keep a short linear transformation
together when splitting it would hide the order \cite[Expose the procedure]{skill:coding}.

\paragraph{Configuration-driven frameworks.} The native configuration is already the
visible execution plan:
\begin{itemize}
  \item keep its stage composition and physics choices; add thin domain adapters only where
        needed;
  \item use native includes or composition before custom preprocessing;
  \item generate repetitive catalogues only from frozen, validated inputs, with deterministic
        regeneration evidence \cite[Ownership and framework composition]{skill:implementation}.
\end{itemize}

\section{Implement a narrow path and inspect its evidence}

\begin{headline}
  Prove the stage structure on one bounded path first. Passing that slice supports only
  its declared coverage.
\end{headline}

In the example, one approved nominal input and one existing observable establish the
stages before expanding to every sample and systematic. The slice validates nothing about
untouched branches \cites[Establish the scientific contract]{skill:coding}[Validation
  order]{skill:implementation}.

\paragraph{Before changing anything, preserve the baseline:} source revision, input
identity, configuration, command, environment and comparison artifacts. After the
rewrite, run the same contract and compare stages, not just the final picture.
Machine-readable counts and bin contents keep exact evidence; plots show patterns and
unexpected shapes \cites[Performance]{skill:implementation}[Validation
  order]{skill:implementation}.

\paragraph{Iterate:}
\begin{enumerate}
  \item Change one coherent stage or interface inside the owned boundary.
  \item Run the affected focused check and the bounded intended input when the stage needs
        real-artifact validation.
  \item Inspect intermediate accounting and output differences against the declared expectation.
  \item Fix the first unexplained discrepancy, record the cause, and rerun the affected checks
        before expanding downstream.
\end{enumerate}
At a coherent integration checkpoint, run the owning required suites. Repeating unchanged
test runs adds no new scientific argument; broaden or repeat validation for new changes,
failures or open concerns \cite[Numerical evidence]{skill:implementation}.

\section{Choose tests that protect the scientific claim}

\begin{headline}
  For personal scripts and small libraries, acceptance is led by running the intended data
  and inspecting the numbers and plots. Tests guard what a plausible plot could hide.
\end{headline}

Keep a small automated suite for invariants and demonstrated bugs; reusable numerical
helpers benefit from a few independently calculated examples, while short recipes usually
need less machinery \cite[Numerical evidence]{skill:implementation}. For the histogram,
Table~\ref{tab:code-tests} lists errors that survive a visually plausible result. Choose
only the cases the contract needs, calculate their expected behavior independently, and
keep the reasoning with the test \cite[Numerical evidence]{skill:implementation}.

\begin{table}[H]
  \small
  \begin{tabularx}{\textwidth}{L{0.36\textwidth}Y}
    \toprule
    Case & Exposes \\
    \midrule
    An event exactly at a cut boundary & The wrong inequality. \\
    \rowcolor{ctpSurface0!30}
    Objects with equal values & An unintended tie rule. \\
    Signed weights & A mistaken normalization. \\
    \rowcolor{ctpSurface0!30}
    A value at or beyond a bin edge & Wrong flow handling. \\
    \bottomrule
  \end{tabularx}
  \caption{Test cases for the histogram example.}
  \label{tab:code-tests}
\end{table}

\review\ A regression test must fail against the incorrect behavior it guards. Copying
the implementation's formula into the expected value repeats its assumption. Tests that
freeze source wording, import lists, internal call structure or an evolving command
interface need a demonstrated requirement; they do not replace scientific evidence
\cite[Numerical evidence]{skill:implementation}.

\section{Find the first mismatch upstream}

\begin{headline}
  Validate in dependency order and inspect the first applicable stage before relying on
  the next. A final
  histogram can agree through compensating errors.
\end{headline}

\begin{figure}[H]
  \centering
  \begin{tikzpicture}[node distance=4mm,
      stage/.style={check,text width=128mm,align=left}]
    \node[stage] (a) {\textbf{1. Source and event accounting}: same intended source,
      file and event coverage, explicit treatment of missing or repeated inputs.};
    \node[stage,below=of a] (b) {\textbf{2. Schema, units, identity}: compatible branches
      and units; object multiplicities and event/object identity make later comparisons
      meaningful.};
    \node[stage,below=of b] (c) {\textbf{3. Weights and normalization}: inspect factors
      before their product and its total; keep signed-weight conventions.};
    \node[stage,below=of c] (d) {\textbf{4. Raw and weighted cutflows}: the first selection
      stage where counts or weighted sums diverge.};
    \node[stage,below=of d] (e) {\textbf{5. Inclusive distributions}: object and event
      behavior before particular regions.};
    \node[stage,below=of e] (f) {\textbf{6. Regions and observables}: overlaps, yields and
      the measured quantities under the declared definitions.};
    \node[stage,below=of f] (g) {\textbf{7. Systematic coverage}: which variations execute,
      which paths stay absent or untested.};
    \node[stage,below=of g] (h) {\textbf{8. Response and closure}, where relevant:
      matched/missed/fake accounting, response orientation, acceptance, efficiency,
      closure, stress cases.};
    \node[result,text width=128mm,align=left,below=of h] (i) {\textbf{9. Final result and
        covariance}: interpret only after the upstream contract is checked.};
    \foreach \x/\y in {a/b,b/c,c/d,d/e,e/f,f/g,g/h,h/i}
    \draw[flow] (\x.south)--(\y.north);
  \end{tikzpicture}
  \caption{Validation order and the question each stage resolves
    \cite[Validation order]{skill:implementation}.}
  \label{fig:validation}
\end{figure}

\paragraph{Example.} Selected event identities, raw counts and observable values agree,
but the weighted histogram differs. That focuses the next look on weight factors,
normalization and filling conventions; it does not say which is wrong, so compare their
intermediate evidence. If input coverage already differed, that earlier discrepancy comes
first. This illustrates the order; it diagnoses no current run.

Record the cause at its owning layer: a formula in its canonical implementation, a host
path in runtime configuration, an intended convention in the confirmed scientific
definition. Keep the failed comparison as provenance, so a later reader sees why the
expectation changed \cites[Ownership and framework composition; Numerical
  evidence]{skill:implementation} [Source traceability and
  corroboration]{skill:report-production}.

\section{Measure performance and finish with a reading guide}

\begin{headline}
  When performance matters, measure before and after on a declared basis; an unexplained
  failure of the agreed gate blocks acceptance even when the numbers agree.
\end{headline}

\begin{itemize}
  \item \textbf{Declare the basis:} input size, host, CPU allocation, I/O location, command,
        wall time and peak resident memory, and the allowed regression before the rewrite.
        A comparison across different storage or CPU states that limitation.
  \item \textbf{Keep what has measured value:} single-pass, bounded-memory, lazy, vectorized
        or compiled behavior, optimized behind domain-named functions with a clear
        reference implementation.
  \item \textbf{Add caching or concurrency} only for an observed bottleneck or a required
        recovery property \cite[Performance]{skill:implementation}.
\end{itemize}

The handback holds the implementation, validation evidence, limits and a reading guide.
For the example, the guide lets the reader:
\begin{enumerate}
  \item find the entry point and follow the stage outputs;
  \item locate every formula and configuration choice;
  \item tell framework wiring from analysis logic;
  \item run the narrow and the complete checks;
  \item add one observable without duplicating an existing definition.
\end{enumerate}
It includes measured performance where required and states untested coverage \cite[Verify
  and hand back]{skill:coding}.

\review\ Acceptance requires the owning checks, current documentation, and
any applicable performance gate. Physics and numerical conclusions also require
a council review and ultimately the human domain owner
\cites[Verify and hand back]{skill:coding}[Rounds]{skill:council}.

\chapter{Produce and Revise a Report}\label{ch:reports}
\section{Define the reader's question and knowledge}

A report is an explanation for a named reader, built around the question that reader
needs to answer. Start by naming that question, the decision or use, the available
evidence, and the report's status. Then record what the reader can be expected to know.
The session's \code{METHODOLOGY_KNOWLEDGE_BASE.md} (or the equivalent in your project)
should connect terminology, mathematics, coursework, literature, and unfamiliar methods
to the planned explanation \cite[Audience and authority; Design the
  argument]{skill:report-production}.

\textbf{Worked example.} Suppose the code rewrite from Chapter~\ref{ch:code}
needs a report answering: ``What evidence supports the claim that this rewrite
preserves the agreed transformation, and what remains untested?'' The input
is its stage contract, source revisions, baseline and candidate artifacts,
comparison records, and performance evidence when required. The report should
explain the acceptance argument; it need not reproduce every execution log.
This example describes a report design, not a claim that the rewrite passed.

\contract\ Assume the reader understands histograms and the analysis context,
but explain any unfamiliar comparison method before using it to accept or reject
agreement. Prefer methods already grounded in the reader's coursework or mapped sources.
A custom diagnostic needs its motivation, definition, limitation, and status
stated explicitly. An unfamiliar method cannot become self-explanatory merely
because it produces a plot \cite[Design the argument]{skill:report-production}.

\section{Inspect the whole artifact and choose its form}

For an existing report, read its complete source, outline, generated evidence, and
rendered PDF before editing. Choose a local patch when the organization still answers the
question; a structural revision when the argument needs consolidation or reordering; and
a rewrite when the organizing model no longer serves the reader. Repetition, appended
findings, and chronology crowding out current conclusions are reasons to reconsider the
structure \cite[Inspect before editing]{skill:report-production}.

Choose the lightest form that gives the reader a clear route. A memo can answer one
bounded comparison; an article can develop one continuous study; a book can carry several
dependent questions and durable reference material. For the illustrative rewrite
comparison, a short article may suffice if there is one argument and limited evidence. A
broader analysis guide may need chapters. Scope and dependency determine the form, not
the template already open \cite[Choose the artifact]{skill:report-production}.

Before drafting, assign each piece of material a role. In our example:
\begin{enumerate}
  \item The opening states the question, bounded coverage, current finding, and review status.
  \item The method explains the frozen contract and how the comparisons test it.
  \item The evidence follows stage order and makes differences inspectable.
  \item Interpretation explains what the evidence supports and what it cannot.
  \item An appendix preserves exhaustive comparisons, commands, and provenance needed to
        reproduce the check.
\end{enumerate}
This is an example outline. The canonical workflow requires a coherent argument,
not these exact headings \cite[Design the argument]{skill:report-production}.

\section{Make every important claim independently checkable}

The bibliography is part of the reader's working apparatus. A reader who can identify a
source, locate the relevant passage, and return to the report's use of it checks the
report far more easily. New substantial reports and substantial rewrites therefore use
shared BibTeX-format \code{.bib} databases, \code{biblatex} with Biber, linked citations,
and bibliography back-references \cite[Bibliography and reader-verifiable
  pinpoints]{skill:report-production}.

Keep the shared library in one directory. External literature has one owner in one file;
maintained project references use a second; dated evidence is grouped by research topic
in further files \cite[Register the source you actually inspected]{skill:library}. Each
source has one stable key and the same descriptive short name across reports. A study
entry may identify a defined evidence collection: the citation then names the particular
member record, while its manifest retains input identities and hashes. A rerun must not
silently replace the version used as evidence by an earlier report.

Three pieces of information do different jobs:
\begin{description}
  \item[Source identity] The bibliography supplies checked author or organization, title,
        edition/version, date, stable identifier where available, and access restriction. A
        repository source instead needs its exact relative path, owning role, and inspected
        revision or date.
  \item[Claim pinpoint] The citation beside the statement names the source's printed page,
        section, equation, table, figure, or record supporting that claim.
  \item[Return route] The bibliography's ``cited on'' pages lead back to this report's uses of
        the source. They are report pages, not source pinpoints.
\end{description}
These distinctions also apply to the report's tables and captions
\cite[Source traceability and corroboration; Bibliography and reader-verifiable pinpoints]{skill:report-production}.

A citation names the source that owns the claim. The hierarchy is:
\begin{itemize}
  \item \textbf{Evidence claims} cite a machine-generated source (runtime output, log,
        executed configuration, code \code{file:line} at a commit) or a human-generated one
        (group documentation, wiki pages, papers, the project's maintained manual). Small cited
        artifacts are copied with their SHA-256 into the session's \code{evidence/} folder,
        because scratch storage is not durable.
  \item \textbf{Session records} (brief, state, result, review and worker records) are cited
        only for the owner's decisions and for review acceptance, never as evidence of a measured
        or implemented fact.
  \item \textbf{Recommendation and deprecation claims} cite the entry in the project's maintained manual or
        register, when one exists. It owns their external citations and records where
        documentation and code disagree. A missing entry is added to the register first.
  \item \textbf{Physics reasoning and tool usage} cite human-facing documentation first: official
        group pages, wiki topics with their revision, tool documentation, and papers. A merge request, commit or \code{file:line} is cited as well
        when that documentation does not point to the change, or when a report explains
        exactly how a tool behaves.
\end{itemize}
Citation and cross-reference links are blue (\code{colorlinks} with \code{ctpBlue}), never
hidden. Reports use descriptive names instead of labels borrowed from other projects
\cite[Source traceability and corroboration; Source hierarchy; Bibliography and
  reader-verifiable pinpoints; Writing character]{skill:report-production}.

\hook{reports-sources}

The following is citation syntax, with placeholder keys and locators to replace after
inspecting the actual sources:
\begin{lstlisting}
\cite[p.~23, sec.~3.9]{source-key}
\cites[pp.~10--11]{first-key}[sec.~2]{second-key}
\end{lstlisting}
Each source in a grouped citation receives its own pinpoint. A bibliography entry's
\code{pages} field usually describes the publication's extent or article number; it does
not identify the passage supporting every claim. For a paper, check the printed
pagination of the cited edition. For Markdown, cite the exact section heading. If a
source is unavailable, mark the missing check rather than guessing a page
\cite[Bibliography and reader-verifiable pinpoints]{skill:report-production}.

In the comparison example, cite the stage record for the agreed transformation, the
precise source revision for implementation, and the comparison artifact for observed
agreement or disagreement. A workflow citation can justify the order of checks; it cannot
substantiate a measured result. A dirty working copy must be identified as such rather
than represented as the published base commit \cite[Source traceability and
  corroboration; Bibliography and reader-verifiable pinpoints]{skill:report-production}.

If a figure includes a source PDF panel, record the caption's printed page and figure
identifier separately from the extraction command's \code{page=} index. Front matter can
make these numbers differ. For derived figures, identify both the input evidence and the
method when they have different owners. Preserve access restrictions when reproducing
material \cite[Choose form by information structure; Bibliography and reader-verifiable
  pinpoints]{skill:report-production}.

\subsection{Keep the editable source readable}

Format authored LaTeX with \code{latexindent} and the project's \code{latexindent.yaml}:
two-space indentation and approximately 90-character prose lines. Preserve paragraphs,
deliberate equation and table breaks, and code listings. Check formatting before handback
using the \code{check_report_sources.py} script of the report-writing skill; its
\code{--write} option applies the formatting. Generated tables should acquire readable
source through their generator, so the next build retains it.

Use \code{siunitx} for numerical results, quantities and units. For example:
\begin{lstlisting}
\SI{25}{\giga\electronvolt}
\num{3084925}
\SI{0.0193}{\percent}
\unit{\pico\barn}
\SIrange{20}{100}{\giga\electronvolt}
\end{lstlisting}
Use \code{\SI} rather than \code{\qty}, which the \code{physics} package uses for
delimiter sizing. Keep the original precision and uncertainty notation. Dataset
identifiers, hashes, software versions, dates and algebraic constants retain their own
notation. Numeric table columns can use \code{S} alignment when it helps comparison. The
quantity scan flags likely omissions; checking the prose, captions and tables still
requires distinguishing values from identifiers \cite[Readable source and numerical
  presentation]{skill:report-production}.

\section{Build the explanation around useful evidence}

Introduce a table or figure before asking the reader to interpret it. State what it
compares, what is held fixed, and which feature answers the question. For the rewrite
example, a compact stage-comparison table can expose the first mismatch; a labelled
distribution comparison can show its shape. Exhaustive per-file logs belong in supporting
evidence unless a particular entry explains the conclusion \cite[Design the argument;
  Choose form by information structure]{skill:report-production}.

For a personal working report the target is a short time to read and digest, not a short
page count; note style is fine while it stays clear. A recommended order of preference
for explaining something, which you may reorder for your readers, is:
\begin{enumerate}
  \item plots, sketches and diagrams; then commented equations with each term explained;
  \item lists and tables;
  \item prose.
\end{enumerate}
Build the aids first, then write the prose that presents the argument and connects them.
A code or pseudo-code snippet may explain what code does when it is clearer than prose. If
the reader would need to re-read a paragraph to follow it, it should become an aid. A diagram
explains a relationship; it does not by itself establish numerical agreement. Keep source
data separate from presentation when generating repeated tables so prose and LaTeX do not
acquire independently maintained copies of the same facts \cite[Choose form by information
  structure]{skill:report-production}.

Three devices from the report template make the structure visible \cite[Headlines,
  choices and colour]{skill:report-production}:
\begin{description}
  \item[Headline box] The section's takeaway in one or two sentences, first thing in the section,
        so the payoff does not wait until its end.
  \item[Choice box] Every choice for the reader, itemized, with a recommendation (blue), an
        alternative (mauve) and links to the supporting plots or tables.
  \item[Colour] Shaded table rows, marked outliers and coloured list markers where they help the
        eye, each colour with one stable meaning and a textual cue.
\end{description}
When the story jumps, reorganize: a definition or equation sits next to its first use.
Once a question is settled, a new dated revision may retire the presentation that served
it; earlier revisions keep the history \cite[Inspect before editing]{skill:report-production}.

Useful visual conventions are portrait pages, readable spacing, explicit reading routes,
restrained semantic color, and legible code blocks. Give every color a stable labelled
meaning. A green ``evidence'' heading can identify an observed record, but a proposed
expectation must not look like an already verified result. Inspect grayscale rendering
when color carries meaning \cite[Colour convention]{skill:report-production}.

\section{Iterate figures and the plotbook before finishing the companion}

Start from the equation. Each plotted quantity is an equation, numbered in the report or
written in the plot's how-to-read note, and the plot shows it as a function or its terms
as histograms, with the equation's symbols in the labels. The standard grammar of the
field comes first: a distribution with extra panels for a ratio, a pull or a diagnostic
such as the stacked terms of a bin decomposition. A new plot type makes the reader learn
a new grammar, so it must earn its place, and its note says why. Every axis is labelled:
a shared $x$ axis once, under the lowest panel, and each panel its own $y$ axis; a
mathematical expression is preferred to words unless a name reads better. If you maintain
your own plotting library, add a missing reusable capability there, with a test, rather
than drawing around it in a recipe \cite[Design and iterate at two scales; Reuse the
  implementation]{skill:plot-production}.

Begin with the reader's task and the evidence needed to perform it. Coverage and page
count are separate: a dense plotbook makes useful comparisons easy without discarding
unique measurements. Recover the applicable binning, statistical conventions and plotting
machinery, then plan candidate groupings. The first plan cannot establish how a
complicated overlay will read before it is rendered \cite[Design and iterate at two
  scales; Reuse the implementation]{skill:plot-production}.

There are two connected revision loops:
\begin{description}
  \item[Figure.] State the question, choose an encoding, render actual data and test whether a
        series can be traced from its legend through the main and ratio panels. Colors, line
        styles and markers are tools. Keep them when they help decoding; revise them when they
        disappear or compete at final size.
  \item[Plotbook.] Check coverage, group comparisons, compose pages and inspect the whole
        sequence. Decide what rows and columns mean. Bring alternatives together when
        simultaneous viewing helps; separate them when they obscure one another. A sparse page
        with tiny plots calls for reconsidering the composition.
\end{description}
A page-level finding can reopen figure design. Neither a fixed grid nor a
historical source/flavour split is a universal solution. Agents normally iterate
within the confirmed contract. An unresolved scientific choice or a conflicting preference
is settled by a recorded default that the owner reviews later, unless no default is
defensible \cite[Design and iterate at two scales]{skill:plot-production}.

A recommended habit is to read the plotbook before the report, for a first impression, so
that the report's interpretation does not anchor the reading (the Einstellung effect).
The plotbook must therefore be readable on its own \cite[What a plotbook holds;
  Established appearance]{skill:plot-production}:
\begin{itemize}
  \item legends and colour keys sit inside the plot, and axis labels name the compared stages,
        for example $(N_\mathrm{BM}-N_\mathrm{BH})/N_\mathrm{BH}$, instead of a title or an
        explanation block atop the page;
  \item a plot whose reading is not obvious carries a short how-to-read note pointing to the
        report section or numbered equation that explains it, and no task-specific term is left
        undefined;
  \item the plotbook holds plots: a long table such as a per-event longtable does not belong, a
        small selected table may;
  \item a page with a large empty area is rejected: the render gate fails a blank band taller
        than \SI{25}{\percent} of the page height unless an exception is recorded;
  \item once results transfer between validation cohorts, the development cohort and one
        transfer-check summary replace side-by-side cohorts.
\end{itemize}

Prove the densest or least certain group in its actual page placement. Audit scale
choices across every affected family, including tails, zero/signed bins and the
uncertainty envelopes actually drawn. Reserve annotation space after overlays and final
scales: a later axis-limit change can undo earlier clearance. Numerical preservation and
geometry checks support this work but cannot judge whether the comparison communicates
its effect \cite[Scales and final geometry; Render and accept]{skill:plot-production}.

For a companion report, establish a provisional argument early but settle the plotbook's
comparisons and page map before final prose. Explain why each comparison answers the
question; assume domain knowledge, not participation in the agent sessions. Introduce a
population before its ratio and use a short worked example where it removes a difficult
conceptual step. Give recurring conventions one home. If writing exposes missing
evidence, return to the plotbook and recheck affected routes rather than writing around
the gap \cite[Write the companion after settling the plotbook]{skill:report-production}.

After the complete draft exists, read it as one argument. Remove repeated context,
reconcile old and new conclusions, move chronology out of the main explanation, and
define terms before use. If the comparison has changed the conclusion, revise the opening
and interpretation together rather than appending a corrective paragraph at the end
\cite[Design the argument; Writing character]{skill:report-production}.

\section{Use separate review tracks and bounded corrections}

A report can be grammatical and wrong, numerically sound and unsupported by its
citations, or correct in source but unreadable in the PDF. Table~\ref{tab:reviews}
separates these checks. Internal review prepares the artifact for the owner's own
inspection; it does not replace their judgment.

\begingroup
\small
\setlength{\LTcapwidth}{\textwidth}
\begin{longtable}{L{0.17\textwidth}L{0.34\textwidth}L{\dimexpr0.49\textwidth-6\tabcolsep}}
  \caption{Report review tracks and the evidence each supplies
    \cite[Source traceability and corroboration; Tiered review]{skill:report-production}.}
  \label{tab:reviews}\\
  \toprule Track & Question & Required evidence \\ \midrule
  \endfirsthead
  \toprule Track & Question & Required evidence \\ \midrule
  \endhead
  Source & Does the cited passage or record support the actual claim? & Checked
  source identity, edition, pinpoint, and claim; unresolved conflicts recorded. \\ \addlinespace
  Copy and structure & Can the reader follow the complete argument in the
  intended voice? & Independent complete-source review for a substantial
  report; findings resolved by the writer. \\ \addlinespace
  Domain / numerical & Are the definitions, calculations, and interpretation
  justified? & Capable-model review of the complete artifact when conclusions
  depend on such claims, plus required human domain review. \\ \addlinespace
  Build and bibliography & Does the source compile with all references and
  citation data resolved? & LaTeX and Biber logs, no unresolved keys or references,
  and the final generated PDF. \\ \addlinespace
  Rendered pages & Are text, tables, figures, labels, and navigation usable? &
  Complete applicable mechanical checks, then bounded independent inspection of
  new layouts, representative families and flagged pages; reuse accepted evidence. \\
  \bottomrule
\end{longtable}
\endgroup

During drafting, review changed claims and their immediate dependencies. Once bounded
findings are resolved, a substantial report needs the complete source corroboration and
copy gates; apply the complete capable-model gate when physics or numerical conclusions
require it. After a local correction, repeat only the affected checks unless it changes
structure, argument order, source hierarchy, global terminology, document-wide
conventions, or claims outside that scope \cite[Source traceability and corroboration;
  Tiered review]{skill:report-production}.

For example, a wrong source page requires rechecking the passage and affected claim, then
rebuilding its citation. A changed normalization can alter the method, tables, plots,
interpretation, and conclusion; that is a wider domain and source correction. A long path
overflowing a cell is a rendering defect, not automatically a scientific defect. Resolve
each finding at the layer that owns it \cite[Production workflow; Tiered
  review]{skill:report-production}.

Complete coverage can combine prior accepted evidence and bounded rechecks; it need not
repeat the same work in every reviewer session. Programmatic checks cover geometry, blank
bands, printed values, inventory and navigation where applicable. Missing checks are not
assessed, and an unblocked lint result alone is not readiness. Before visual dispatch,
validate the explicit model and page selection with the render gate's preflight. Use an
economical image-capable model for remaining reading tasks. Repagination reopens
navigation and composition; it does not automatically invalidate accepted page-body
readability. Expensive or exhaustive review requires a concrete unresolved problem
\cite[Tiered review]{skill:report-production}.

Give a copy reviewer the complete editable source. For visual delegation, rasterize first
and provide relevant labelled PNG or JPEG contact sheets and selected pages in a scratch
directory; never hand the visual reviewer a PDF. A vision report is evidence about
appearance and needs comparison with the source and intended rendering before it becomes
a claim about correctness \cite[Tiered review]{skill:report-production}.

Give the independent reviewer the reader brief and artifact before the author's verdict.
Ask it to reconstruct the question, findings and supporting comparisons without the
session narrative. It may find a structural problem even when the sentences are correct.
On rendered pages, ask it to trace a series, read the labels and identify the intended
effect. Contact sheets establish overview coverage; individual pages at reading size are
needed for decoding and small text. An unreadable review image means that feature is not
assessed \cite[Review packet; Figure and page tasks; Companion comprehension
  tasks]{skill:visual-review}.

Record the exact artifact, page, observation, reader impact and recheck result. Verify
alleged defects against the image: agreement with criticism does not prove inspection.
Iterate until reading tasks work and material defects are resolved; neither a fixed
number of passes nor repeated cosmetic changes establish quality \cite[Findings and
  revision]{skill:visual-review}.

\section{Build, inspect, and deliver the exact PDF}

The recommended bibliography setup uses \code{backend=biber}, \code{style=alphabetic},
\code{sorting=nyt}, and \code{backref=true}. Use readable, unambiguous short source
labels where they help. A manual build sequence is shown below; \code{latexmk} runs the
same steps and repeats them until the references settle:
\begin{lstlisting}
pdflatex -interaction=nonstopmode -halt-on-error report.tex
biber report
pdflatex -interaction=nonstopmode -halt-on-error report.tex
pdflatex -interaction=nonstopmode -halt-on-error report.tex
\end{lstlisting}
Inspect both logs for missing or duplicate keys, malformed data, unresolved references,
rerun requests, and material layout warnings. Check the rendered citation-to-entry links
and bibliography return links as well as page appearance \cite[Bibliography and
  reader-verifiable pinpoints]{skill:report-production}.

The authored \code{.tex} and \code{.bib} are the reconstructible source. The PDF and
ordinary bibliography/build intermediates remain generated. A clean build proves that the
source can be rendered; the source, copy, domain, and visual reviews establish the
additional evidence needed for delivery \cite[Acceptance checklist; Bibliography and
  reader-verifiable pinpoints]{skill:report-production}.

Hand back the exact PDF and source paths, review evidence, and remaining limits. For a
requested plotbook companion, deliver the matched pair with checked routes; raw logs
remain supporting material. Keep reviewable copies in one collection, delivered as soon
as they are ready: a same-day revision replaces its file, a later date goes alongside,
with every revision's hashes and human-review state in a manifest \cite[Review and
  release]{skill:report-production}. \hook{reports-delivery} Record authorized external
delivery, notification, and source publication separately \cite[Review and
  release]{skill:report-production}.


\chapter{Run, Monitor, and Diagnose Remote Work}\label{ch:remote}
\section{Preflight the actual execution boundary}

Use this flow when validation needs data, software, or resources outside the current
writing environment, or when a job outlives an interactive step. Start with the owning
repository's host and execution documentation. Establish the literal host, exact target
and owner, capability, and consequence. Synthetic checks can run on the writing host;
service-dependent checks need the documented capable environment \cite[Remote execution
  and interruption]{skill:remote}.

Access to remote data does not by itself require remote execution. Choose the execution
host for its runtime, input access and capacity, then test the exact input from that
host. Before relying on a route, distinguish a verified read from an unverified
submission path: a successful storage read does not validate batch submission.
\hook{remote-hosts}

Probe the chosen SSH alias non-interactively before dependent work. With a shared control
master, issue one bounded command at a time through that alias. After a timeout, broken
pipe, multiplexing error, or unexpected disconnect, stop issuing calls and inspect the
master once. Do not replace the user's transport as an incidental repair. If
authentication is absent, the user establishes one interactive login; the agent never
requests a password or certificate passphrase \cite[SSH control masters]{skill:remote}.

The authorization must name the intended remote action. Inspecting logs is different from
submitting or cancelling a job. A bounded pilot may be a useful first execution when its
submission is authorized; use it to check the input, configuration, output schema, and
actual resource behavior before expanding the run. Keep unbounded analysis inside the
scheduler rather than on a login node. Source edits stay in the explicitly owned local
checkout or worktree, never through a mounted remote filesystem \cite[Remote checkouts
  and execution; Scheduler jobs]{skill:remote}.

\section{Preserve the parent contract before waiting}

Record exact job IDs, status surfaces, evidence paths, allowed and forbidden actions,
monitoring interval and deadline, terminal validation, and downstream artifacts. Include
separately any report, review, source publication, delivery, and notification
requirements. This lets the parent recover the entire task when a wait returns rather
than treating scheduler completion as the endpoint \cite[Durable completion
  contract]{skill:jobs}.

Prefer an existing authorized programmatic monitor. It can observe exact external tasks
and write local snapshots without spending model calls on polling. The documented monitor
preserves atomic current state, append-only observations, a terminal result, and an
exclusive lock. Snapshot expiry matters: a stale record is missing current evidence, not
proof of a successful task \cite[Monitor selection; Programmatic job
  monitors]{skill:jobs}.

Where the client provides one, a parent wait call binds a reviewed contract and monitor
hashes, a finite duration, a handback policy, and any exact notification authorization.
It returns evidence to the parent. A client that lacks such a tool must verify its own
integration and parent identity handling before relying on that behavior. Otherwise
record the gap and use the documented bounded fallback or explicit manual recovery of
existing evidence \cite[Client scope; Parent wait contract]{skill:jobs}.

Before unattended work, preflight the exact external read and write roots, permission
configuration, and ownership. A permission prompt is a handback condition; widening paths
or using blanket auto-approval does not solve the contract. The specific configuration
rules belong to each client, while the requirement to respect the task boundary applies
across clients \cite[Permission and access preflight; Client scope]{skill:jobs}.

\section{Interpret the handback before changing anything}

A monitor supplies observations. Table~\ref{tab:handback} describes the next reasoning
step for common handbacks; repair or retry requires its own recorded authorization.

\begingroup
\small
\begin{longtable}{L{0.27\textwidth}L{0.66\textwidth}}
  \caption{Remote handback and the bounded next inspection
    \cite[Programmatic job monitors; Terminal handling]{skill:jobs}.}
  \label{tab:handback}\\
  \toprule Observation & Next inspection \\ \midrule
  \endfirsthead
  \toprule Observation & Next inspection \\ \midrule
  \endhead
  Active, fresh snapshot & Continue the contracted wait within its lease. \\
  Stale or absent snapshot & Establish monitor liveness and current evidence;
  do not infer external job success. \\
  Terminal clean & Validate the predeclared files, events, schema, checksums,
  provenance, normalization, and storage before downstream use. \\
  Terminal with failures & Inspect the exact failing task/index and relevant
  logs; preserve successful outputs and failed-attempt evidence. \\
  Owner/scope/invariant drift & Reconcile the exact job identities, specification,
  and expected counts; stop expansion until the discrepancy is understood. \\
  Transport failure & Apply the shared SSH stop-and-inspect rule; distinguish
  loss of observation from failure of the analysis. \\
  \bottomrule
\end{longtable}
\endgroup

Suppose an array reports mixed completion. First identify the missing or failed indices
and their exit evidence, then check which outputs exist and which input coverage they
represent. Preserve successful evidence. The smallest next decision may concern only the
failed subset, but monitoring permission alone does not authorize its resubmission.
Record the cause, bounded proposal, and required authorization instead of restarting the
entire production by assumption \cite[Terminal handling]{skill:jobs}.

If the scheduler reports success but event coverage is incomplete, downstream analysis
remains held. Follow the upstream validation order in Chapter~\ref{ch:code}: input and
accounting before distributions and final results. Logs, schema inspection, and exact job
evidence should drive the diagnosis. Add focused logging only when the current evidence
cannot distinguish the plausible causes \cite[Validation order]{skill:implementation}.

\section{Handle client loss and notification uncertainty}

A documented wait separates observation state from delivery state. A detached notifier
can continue observing existing monitors after client loss and record an orphaned result;
it never starts an agent or model. An explicitly selected recovery owner must reconcile
the result and the remaining contract \cite[Client loss and notification; Compaction and
  recovery]{skill:jobs}.

Email is disabled unless the exact command and destination are authorized. An interrupted
attempt can leave delivery externally uncertain. Do not retry an ambiguous attempt
automatically: reconcile the notifier and transport evidence first. Acknowledging a wait
means its evidence was consumed; it does not validate the output or complete the
scientific task \cite[Client loss and notification]{skill:jobs}.

At handback, update the parent state with terminal checks, known failures, remaining
analysis or report work, and one next safe action. Restore the same obligations after
compaction. A job leaving the queue, a worker returning, a PDF compiling, and a local
email command succeeding each support a particular claim; none alone closes the whole
parent contract \cite[Completion evidence]{skill:jobs}.

\chapter{Coordinate Changes and Hand Them Off}\label{ch:handoff}
\section{One writer per boundary}

\begin{headline}
  Use coordination when work crosses repositories, another agent contributes, or ownership
  moves between machines or clients. Each boundary has one writer and a precise packet.
\end{headline}

Start by naming the packet's fields \cite[Dispatch packet; Review]{skill:delegation}:
\begin{itemize}
  \item the deliverable and the evidence;
  \item the writer and its branch or worktree;
  \item forbidden actions;
  \item verification and the handback condition.
\end{itemize}
A precise packet lets a worker finish a bounded contribution without silently inheriting
the parent objective.

\paragraph{Repositories and worktrees} \cite[Start a writing task; Commit and hand
  off]{skill:git}.
\begin{itemize}
  \item Software repositories retain independent histories. A metadata or project-management
        repository owns policy, session records and approved report sources; it is not a
        container for their code history.
  \item Parallel software changes and overlapping report edits use separate worktrees.
  \item A shared checkout can coordinate non-overlapping state paths under one integrator.
  \item Two tasks needing the same report or policy file require serialization or one explicit
        writing owner.
\end{itemize}

\begin{figure}[H]
  \centering
  \begin{tikzpicture}[node distance=6mm and 10mm]
    \node[work,text width=40mm] (c) {Document the input/output contract between the
      layers};
    \node[work,text width=40mm,right=of c] (w) {Each worker gets one exact repository and
      interface};
    \node[check,text width=40mm,right=of w] (i) {The integrator checks the combined evidence
      in dependency order};
    \draw[flow] (c)--(w);
    \draw[flow] (w)--(i);
  \end{tikzpicture}
  \caption{A cross-repository change. One implementation owner owns each physics formula and
    convention. A locally passing helper does not demonstrate that its consumer receives the
    expected units, schema or normalization \cites[Ownership and framework composition;
      Validation order]{skill:implementation}[Dispatch packet; Review]{skill:delegation}.}
  \label{fig:handoff-layers}
\end{figure}

\section{Review a contribution before accepting it}

\begin{headline}
  Review the handback's evidence, not its confidence. Physics and numerical conclusions go
  to a council.
\end{headline}

A useful worker handback contains \cite[Choosing a role; Review]{skill:delegation}:
\begin{itemize}
  \item the actual diff and the checks run;
  \item evidence paths and remaining limits;
  \item a reading route.
\end{itemize}
Independent copy review is read-only feedback to the report writer; domain and
rendered-page review remain separate.

\begin{rulebox}[Mechanical workers]
  A mechanical worker does not decide ambiguous ownership, job control, destructive cleanup,
  or final physics conclusions.
\end{rulebox}

\paragraph{The parent then} \cites[Review]{skill:delegation}[Numerical
  evidence]{skill:implementation}:
\begin{enumerate}
  \item compares the contribution with the original packet and the remaining contract;
  \item resolves findings at their owning layer;
  \item re-runs the affected checks in the relevant environment;
  \item performs the required integration checks once the pieces are coherent;
  \item records deferred or unavailable checks explicitly: they limit the claim the handback can
        support.
\end{enumerate}

\begin{figure}[H]
  \centering
  \begin{tikzpicture}[node distance=5mm and 12mm]
    \node[work,text width=30mm] (brief) {Known fixes first, then one brief};
    \node[check,text width=34mm,right=of brief,yshift=7mm] (m1) {Strong model,\\provider 1};
    \node[check,text width=34mm,right=of brief,yshift=-7mm] (m2) {Strong model,\\provider $n$};
    \node[result,text width=36mm,right=of m1,yshift=-7mm] (v) {Verdicts and findings;
      record who reviewed in which round};
    \draw[flow] (brief.east)--(m1.west);
    \draw[flow] (brief.east)--(m2.west);
    \draw[flow] (m1.east)--(v.west);
    \draw[flow] (m2.east)--(v.west);
    \draw[back] (v.south) to[bend left=30] node[below,lbl]{fix, then round 2}
    (brief.south);
  \end{tikzpicture}
  \caption{The council: one strong model from each provider you can reach reviews the same
    brief in parallel, so a blind spot shared by one model family is less likely to pass
    unseen \cite[Members; Rounds; Record]{skill:council}.}
  \label{fig:handoff-council}
\end{figure}

After two rounds, a finding still open is fixed, or its claim is withdrawn or marked
unsupported. Only a disagreement of judgment that no run can settle becomes a recorded
default or a question to you, and gated actions keep their gates.

\section{Make source durable at the right checkpoint}

\begin{headline}
  Git is the normal source-handoff mechanism across machines. Unique edits are committed
  and pushed before they are left unattended.
\end{headline}

\paragraph{Updating.} Inspect status, branch, worktrees and remotes before an update; pulls
are manual and safe updates fast-forward. Unknown divergence or dirty ownership is
resolved without resetting or overwriting another writer's information \cite[Update
  safely]{skill:git}.

\paragraph{Publishing,} when publication is authorized and classification permits the
destination \cite[Commit and hand off; Material you must not publish]{skill:git}:
\begin{enumerate}
  \item stage only owned paths;
  \item inspect the staged diff;
  \item record checks and limits;
  \item commit coherently;
  \item push to the approved writable remote;
  \item verify the remote commit.
\end{enumerate}
Framework projects have their own fork and contribution rules; consult the repository
instructions before choosing the target.

Scratch is useful for explicitly owned framework checkouts and task worktrees, but it is
not durable merely because it is fast or has space.

\begin{rulebox}[Scratch storage]
  Commit and push unique software edits
  there before leaving them unattended or handing them off. If publication is not authorized
  or fails, record the exact dirty paths and recovery gap in state, and keep the work from
  existing only on scratch \cite[Commit and hand off]{skill:git}.
\end{rulebox}

\begin{table}[H]
  \small
  \begin{tabularx}{\textwidth}{L{0.3\textwidth}Y}
    \toprule
    State & Does not imply \\
    \midrule
    Source checkpoint published & That human report decisions are closed: an authorized
    checkpoint can preserve work while a PDF remains under review. \\
    \rowcolor{ctpSurface0!30}
    PDF delivered & That its source is published. \\
    \bottomrule
  \end{tabularx}
  \caption{Record both states, so the receiving owner knows which evidence is available
    \cites[Commit and hand off]{skill:git}[Review and release]{skill:report-production}.}
  \label{tab:handoff-states}
\end{table}

\section{Give the receiving owner one safe continuation}

\begin{headline}
  Stop the outgoing writer at a coherent checkpoint and hand over one next safe action.
\end{headline}

Give the receiver \cite[Routine continuation and client changes]{skill:recovery}:
\begin{itemize}
  \item the chosen session and exact source state;
  \item current artifacts, active jobs or waits;
  \item remaining checks and the authorization boundary;
  \item the next safe action.
\end{itemize}
The receiver reconciles that record with current Git and runtime evidence before
continuing; Chapter~\ref{ch:daily} distinguishes another display from another writer.

\paragraph{\code{HANDOFF.json}} is useful when a programmatic dependency or cross-harness
recovery needs it \cite[Handoff marker]{skill:jobs}:
\begin{itemize}
  \item it mirrors the Markdown checkpoint atomically;
  \item missing, malformed or held releases do not mean ready;
  \item for a report awaiting human review, the machine status remains \code{in_progress} with
        the report release held.
\end{itemize}

\begin{rulebox}[No speculative release]
  Do not create a release merely because the next step is expected to pass.
\end{rulebox}

\chapter{Recover and Preserve Work}\label{ch:recovery}
\section{Recover context, then verify reality}

\begin{headline}
  An interruption can remove the visible conversation while source, remote jobs and
  monitors survive. Recovered context is a lead; current evidence decides.
\end{headline}

Recover the selected harness's native context when supported, then independently
reconcile the chosen brief and state, the owning Git/worktree, the live agent identity
and the relevant remote status. Restore the whole task contract, including downstream
review and delivery obligations \cite[Recover harness context, then verify
  it]{skill:recovery}.

\paragraph{After a reboot, crash or unexplained disconnect:}
\begin{enumerate}
  \item Identify the actual host and time; use a non-interactive preflight for any required
        remote alias.
  \item Inspect host and storage health without cleanup, remounting or service changes.
  \item Locate the selected agent and its native identity using the client or terminal
        multiplexer that owns it.
  \item Read the session records and reconcile them with Git and current job or monitor evidence.
  \item Choose one desired session to resume, record its owner and checkpoint, and continue only
        within its surviving authorization.
\end{enumerate}

\begin{rulebox}[Stop on]
  \begin{itemize}[leftmargin=1.2em]
    \item an unexpected host;
    \item unknown dirty ownership;
    \item disagreement between state and Git;
    \item required credential renewal;
    \item a recovery action that would overwrite information.
  \end{itemize}
  A restart or reconnect does not automatically resume an old autonomous prompt
  \cite[Startup and recovery, steps 1--8; Stop conditions]{skill:recovery}.
\end{rulebox}

\begin{table}[H]
  \small
  \begin{tabularx}{\textwidth}{L{0.26\textwidth}Y}
    \toprule
    Incident & Response \\
    \midrule
    Harness process: a competing or stuck agent process & Identify the exact processes; a
    bounded repair when authorized. This never authorizes replacing the shared SSH
    master. \\
    \rowcolor{ctpSurface0!30}
    SSH transport & The stop-and-inspect rule of Chapter~\ref{ch:remote}. \\
    \bottomrule
  \end{tabularx}
  \caption{Two kinds of incident. Use incident-specific evidence rather than broad
    process-kill patterns \cite[SSH control masters]{skill:remote}.}
  \label{tab:recovery-incidents}
\end{table}

\hook{recovery-clients}

\section{Preserve the result before reducing session scaffolding}

\begin{headline}
  Once operational work settles, promote lasting evidence into an outcome record that lets
  a later reader recover the result without replaying every step.
\end{headline}

The outcome record holds \cite[When]{skill:lifecycle}:
\begin{itemize}
  \item the objective, result and decisions;
  \item report lineage, commits, jobs, reviews and publication;
  \item unresolved limits, with the leftovers it still depends on.
\end{itemize}
Retain unique evidence until a review confirms that consolidation loses nothing essential.

\begin{rulebox}[Supersession]
  Similar report names or dates do not prove supersession: the newer artifact must
  explicitly incorporate or reference the older evidence. A report awaiting human review is
  still active; its source and unique evidence stay protected until the owner's disposition
  is recorded, whether or not the last agent stopped \cite[Review and
    release]{skill:report-production}.
\end{rulebox}

\section{Choose the recovery source for each kind of content}

\begin{headline}
  Git remotes, complementary copies and rebuilds cover different content. Check each
  separately (Table~\ref{tab:recovery}).
\end{headline}

\begin{table}[H]
  \small
  \begin{tabularx}{\textwidth}{L{0.32\textwidth}Y}
    \toprule Content & Recovery evidence \\
    \midrule
    Approved source and records & Reviewed commit actually present on the approved
    writable remote; a local commit alone is not an off-device copy. \\
    \rowcolor{ctpSurface0!30}
    Eligible local-only reading material & Verified complementary copy or documented
    recovery source; respect destination-specific restrictions. \\
    Generated PDF, plots, build state & Preserved inputs, authored source, and a
    reproducible build or generation command. \\
    \rowcolor{ctpSurface0!30}
    Credentials & Excluded from backup; user follows the authority's reissue or
    initialization procedure. \\
    \bottomrule
  \end{tabularx}
  \caption{Recovery classes: content and the evidence that proves it can be restored.}
  \label{tab:recovery}
\end{table}

\begin{choice}{copy tool for complementary copies}
  \begin{itemize}
    \item \recommended\ A tool that uploads additions and replacements without deleting
          remote-only files: the safer routine choice.
    \item \alternative\ A mirror that can delete, only with specific intent.
  \end{itemize}
  Either way, a copy is not Git history and does not provide the same versioned recovery.
\end{choice}

\begin{rulebox}[Restricted material]
  Neither a private Git approval nor a convenient backup route authorizes sending
  restricted material to another service. Verify destination exclusions before copying.
\end{rulebox}

\hook{recovery-backup}

\section{Review cleanup as a separate decision}

\begin{headline}
  An audit finds candidates; it does not authorize removing them. Reversible removals may
  run without per-item approval; everything else waits for the approver.
\end{headline}

A candidate is eligible when it has \cite[Autonomous with
\texttt{-{}-apply}]{skill:lifecycle}:
\begin{itemize}
  \item known ownership and inactivity evidence, and a current inactive status;
  \item no active dependencies;
  \item required publication and a recovery source;
  \item report lineage preserving unique evidence.
\end{itemize}
Age adds a condition, not a substitute: the lifecycle procedure's default minimum is 30
days. A waiver names exact targets and cannot clear a failed semantic gate.

\begin{figure}[H]
  \centering
  \begin{tikzpicture}[node distance=8mm and 14mm]
    \node[work,text width=44mm] (obs) {Expected topology and a dated observation};
    \node[check,text width=44mm,right=of obs] (elig) {Semantic eligibility per candidate};
    \node[work,text width=44mm,below=of obs] (rev) {Reversible: remove without per-item
      approval};
    \node[open,text width=44mm,below=of elig] (appr) {Otherwise: listed for the approver,
      who returns the report hash};
    \node[check,text width=60mm,below=10mm of $(rev.south)!0.5!(appr.south)$] (act)
    {Reclassify each target immediately before acting; refuse drift};
    \draw[flow] (obs)--(elig);
    \draw[flow] (elig.south west)--(rev.north east);
    \draw[flow] (elig)--(appr);
    \draw[flow] (rev.south)--(rev.south |- act.north);
    \draw[flow] (appr.south)--(appr.south |- act.north);
  \end{tikzpicture}
  \caption{Cleanup. This preserves the ability to explain why a target is unused and how its
    important content could be recovered \cite[When; Command]{skill:lifecycle}.}
  \label{fig:recovery-cleanup}
\end{figure}

\paragraph{Holds.} Dirty work, unknown ownership, active references, missing publication,
uncertain lineage, unavailable remote evidence and target drift. Record the specific
missing evidence rather than turning a broad ``looks old'' judgment into approval. A PDF
being ignored by Git does not make it unimportant during review or handoff.

A reversible removal is, for example, a clean worktree whose branch is already on the
remote. Git worktrees are removed through Git, not raw tree deletion \cite[Approval;
  Restore]{skill:lifecycle}.

\begin{rulebox}[No implied authorization]
  Having this handbook, an audit or a dry run authorizes no cleanup action.
\end{rulebox}

\chapter{Field Checklists}
\section{Before work}
\begin{enumerate}
  \item State the question, bounded outcome, and artifact to inspect.
  \item Identify the source owner, runtime, worktree, and any execution target.
  \item Preserve unknown dirty state; assign one writing boundary.
  \item Use a confirmed decision record when the planning gate applies.
  \item Create the required durable records for substantial or recovery-worthy work; add a
        machine marker only when needed.
  \item Name the checks, applicable reviews, allowed actions, and next stop.
\end{enumerate}
This is the entry contract from Chapters~\ref{ch:daily} and~\ref{ch:decisions}
\cite[rules 1--3]{skill:coordination}.

\section{Before accepting code or a report}
\begin{enumerate}
  \item Check the agreed input and definition before judging downstream output.
  \item For code, inspect stage evidence, relevant independent examples, intended-artifact
        execution, required performance, and the reading guide.
  \item For a report, check the complete argument, source identity and pinpoints, bibliography
        links, and the required independent review tracks.
  \item For reports and other rendered deliverables, build the actual source and inspect the
        actual rendered pages.
  \item State coverage and unresolved gaps; do not promote a bounded check into an unqualified
        scientific conclusion.
\end{enumerate}
Use the detailed gates in Chapters~\ref{ch:code} and~\ref{ch:reports}
\cites[Verify and hand back]{skill:coding}[Acceptance checklist]{skill:report-production}.

\section{Before handoff or interruption}
\begin{enumerate}
  \item Record the owner, exact source state, current artifacts, and next safe action.
  \item Reconcile active jobs, monitors, leases, and uncertain deliveries.
  \item Preserve unique work and verify any authorized publication.
  \item Distinguish attachment to the same agent from transfer to another writer.
  \item Keep a substantial report awaiting human review until the owner's disposition.
\end{enumerate}
These checks preserve the continuation described in Chapters~\ref{ch:handoff}
and~\ref{ch:recovery}
\cites[Commit and hand off]{skill:git}[Recover harness context, then verify it]{skill:recovery}.

\chapter{Worked Example: From a Question to a Reviewed Report}\label{ch:example}

This chapter walks one small task from question to figure to report to review, using the
synthetic counting example of the repository \cite[Two renderers, one evidence
  contract]{example:custom-plotting}. It uses no measured data, so the task is tiny on
purpose: what matters is the order of the steps and the evidence each one leaves. The
same chapter is also how this handbook was built: it follows the report-writing
procedure, and its own review is the last step below.

\section{State the question and the contract}

\begin{headline}
  Write the stage record before drawing anything (Chapter~\ref{ch:code}).
\end{headline}

\textbf{Question.} How many of six synthetic entries fall into each of three fixed bins, and
does the plotted histogram show the same numbers as an independent hand count?

\begin{table}[H]
  \small
  \begin{tabularx}{\textwidth}{L{0.27\textwidth}Y}
    \toprule Field & Record before editing \\
    \midrule
    Authoritative source & The six literal samples in the recipe and the three bin edges
    $0,1,2,3$. \\
    \rowcolor{ctpSurface0!30}
    Named stages & Declare samples and edges; count per bin; compare with the hand count;
    draw; record the printed values. \\
    Oracle & A hand count written down independently of the code: $[2,3,1]$. \\
    \rowcolor{ctpSurface0!30}
    Acceptance evidence & Counts equal the oracle; the rendered page carries the axis
    labels and the printed bin values; the render gate has no blocking finding. \\
    Non-goals & No measured data, uncertainty, fit or physics selection. \\
    \bottomrule
  \end{tabularx}
  \caption{The stage record of the example, written before the recipe is run
    \cite[Establish the scientific contract]{skill:coding}.}
  \label{tab:example-stage}
\end{table}

The plotted quantity is an equation first. With samples $x_j$ and edges
$e_0<e_1<e_2<e_3$, the content of bin $i$ is
\begin{equation}
  N_i = \sum_{j=1}^{6} \mathbf{1}\!\left[e_{i-1}\le x_j<e_i\right],
  \qquad i=1,2,3,
  \label{eq:counts}
\end{equation}
where $\mathbf{1}[\cdot]$ is 1 when sample $x_j$ lies in bin $i$ and 0 otherwise. The figure
shows $N_i$ as a function of the coordinate, with $N_i$ on its vertical axis \cite[Design
  and iterate at two scales]{skill:plot-production}.

\section{Run the recipe and keep its evidence}

\begin{headline}
  The recipe checks its own counts against the oracle before it draws, and leaves its
  evidence next to the figure.
\end{headline}

From the repository root, with a Python environment that has NumPy and Matplotlib:
\begin{lstlisting}
python examples/custom-plotting/render.py --renderer matplotlib \
  --out docs/guide/build/custom
\end{lstlisting}

\begin{figure}[H]
  \centering
  \begin{tikzpicture}[node distance=6mm and 9mm]
    \node[work,text width=24mm] (in) {Samples and edges};
    \node[work,text width=24mm,right=of in] (n) {Count per bin};
    \node[check,text width=32mm,right=of n] (o) {Equal to $[2,3,1]$, sum equal to the
      number of samples?};
    \node[result,text width=24mm,right=of o] (draw) {Draw and write evidence};
    \node[box,draw=ctpText!60,text width=24mm,below=7mm of o] (stop) {Stop};
    \draw[flow] (in)--(n);
    \draw[flow] (n)--(o);
    \draw[flow] (o)--node[above,lbl]{yes}(draw);
    \draw[flow] (o)--node[right,lbl]{no}(stop);
  \end{tikzpicture}
  \caption{What the recipe does.}
  \label{fig:example-recipe}
\end{figure}

The evidence written next to the figure:
\begin{itemize}
  \item the source samples and the formatted values;
  \item the renderer identity;
  \item an inventory of the rendered page;
  \item an applicability record for the render gate.
\end{itemize}
The recipe refuses to reuse an output directory, so old evidence cannot be mixed with new.
Figure~\ref{fig:example} is its output.

\begin{figure}[H]
  \centering
  \includegraphics[width=0.8\textwidth]{build/custom/example-counts.pdf}
  \caption{Synthetic example: entries $N_i$ per bin, as in
    Equation~\ref{eq:counts}. Read the three printed bin values, 2, 3 and 1, and the total
    of 6 at the top right: they agree with the independent hand count. There are no
    measured data.}
  \label{fig:example}
\end{figure}

\evidence\ The oracle and the computed counts agree: $N=(2,3,1)$, total $6$. This supports
only the declared coverage, which is these six samples and three bins
\cite[Numerical evidence]{skill:implementation}.

\section{Write the report around the figure}

\begin{headline}
  A report for this question can be one page, and its claim is limited to the synthetic
  input.
\end{headline}

The one-page report \cite[Source traceability and corroboration; Design the
  argument]{skill:report-production}:
\begin{itemize}
  \item opens with the question, the answer and its status;
  \item introduces Figure~\ref{fig:example} before interpreting it, states what is held fixed
        (samples and edges) and which feature answers the question (the three printed bin
        values);
  \item cites the source that owns each claim: the recipe's evidence for the counts, the stage
        record for the contract, a workflow document only for why the checks come in this order.
\end{itemize}

\gap\ A report that said ``the method works'' would claim more than a synthetic check can
show. The claim that survives is narrower: the recipe's counts match a hand count for these
inputs, and the page carries the values it should.

\section{Review in separate tracks}

\begin{headline}
  Each review track of Chapter~\ref{ch:reports} answers one question. For this handbook
  they run as follows, from the repository root.
\end{headline}

\begin{description}
  \item[Source] \code{skills/report-writing/scripts/check_report_sources.py} checks the chapter
        sources for duplicate bibliography keys, \texttt{\textbackslash addbibresource} targets
        outside the shared library, hand-written bibliography blocks and formatting drift. A
        reader verifies pinpoints by following each citation to its entry.
  \item[Build and bibliography] \code{latexmk} runs \code{pdflatex} and \code{biber} until
        references settle; the \code{.log} and \code{.blg} files are read for unresolved keys or
        references.
  \item[Rendered pages] \code{skills/render-review/scripts/render_gate.py gate} binds the PDF,
        its log and an applicability contract by SHA-256 and writes \code{gate.json} and a review
        packet. The contract records, for each optional check, why it does not apply here. A
        \code{fail} blocks visual review until it is fixed or waived \cite[Command;
          Rules]{skill:render-workflow}.
  \item[Reading tasks] After the gate, a reviewer receives selected page images, never the PDF,
        with the tasks to perform: read the axis labels in Figure~\ref{fig:example}, find the bin
        with three entries, and say what the page claims \cite[Figure and page
          tasks]{skill:visual-review}.
\end{description}

\review\ The gate's status is mechanical evidence about layout, fonts, build errors and
printed values. It is not a verdict that the page reads well; the reading tasks supply that.
The commands and the location of the gate output are in the \code{README.md} beside this
handbook's sources.

\chapter{Make It Yours}\label{ch:yours}

These skills are one person's working style, built during a physics thesis. They are
meant to be taken apart. This chapter says how to adapt them without losing what makes
them useful: rules that live in one place and reach every agent you use.

\section{Start small}

You do not need to use every skill on the first day. Install the whole \code{skills/}
folder, because the skills link to each other, but begin by calling only three:
\code{scientific-mode}, \code{grilling} and \code{unslop}. The first gives each task a
playbook, the second settles decisions before work begins, the third keeps the prose
readable. Work plainly for a while, watch where agents fail, and add a skill or a check
when the same failure shows up twice \cite[Make it yours]{pstack:guide}.

\section{Keep your settings in an overlay}

Do not edit a shared skill to record where your files are or which approver signs a
cleanup. Keep those facts in one personal file next to your project's agent instructions,
for example \code{PERSONAL.md}, and let the skills refer to ``the user'' and ``the
project''. Table~\ref{tab:yours-owners} shows where each kind of rule belongs.

\begin{table}[htbp]
  \small
  \begin{tabularx}{\textwidth}{L{0.3\textwidth}Y}
    \toprule
    Kind of rule & Where it lives \\
    \midrule
    How anyone should work & A shared skill; change it upstream so everyone gets it. \\
    How you work here & The personal overlay: paths, script calls, approvers, delivery. \\
    What the analysis agreed & The analysis repository or its notes: definitions,
    conventions, supervisor decisions. \\
    What happened this week & Session records: \code{BRIEF.md}, \code{STATE.md},
    \code{RESULT.md}. \\
    \bottomrule
  \end{tabularx}
  \caption{Where a rule belongs. A client's own memory feature is read by that client only,
    so it should hold transient context, not rules.}
  \label{tab:yours-owners}
\end{table}

\section{Let lessons reach the file that owns them}

After a task that taught something, run \code{reflect}: it turns the session's lessons
into edits of the skill, overlay or analysis file that owns each rule \cite[Extract
  lessons; Apply]{skill:reflect}. When you correct the same mistake a second time, run
\code{correct}. It prefers a structural fix or a check whose error names the fix over
another sentence of instructions, and keeps a table that pairs each rule with what
enforces it \cite[Fix each class at the highest level that works; Keep the rule
  table]{skill:correct}. A rule that nothing enforces is the first to be forgotten.

\section{Choose how much to delegate}

The default here is to decide, record and report (Chapter~\ref{ch:decisions}): agents
proceed on reversible work and on defensible scientific defaults, and every default stays
visible until you review it. Which actions count as reversible depends on your project.
If your branches are shared, or your review inbox is a shared folder, say so in your
overlay and the gates move accordingly \cite[Decide, record, report]{skill:actions}.

\section{Stay current without being dragged along}

You can install the repository as a plugin and follow it, or copy the skills you want and
own them. If you copy, record the commit you copied from. This repository does the same
with the skills it adapted from others: each records its upstream commit, and a weekly
job opens an issue when an upstream changes. Adopting a change is then a reviewed edit,
not an automatic update.


\backmatter
\printbibliography[heading=bibintoc,title={Sources}]
\end{document}
